\documentclass[12pt]{article}

\usepackage{sbc-template}
\usepackage{graphicx,url}
\usepackage[utf8]{inputenc}
\usepackage[brazilian]{babel}
\usepackage{booktabs}
\usepackage{array}
\usepackage{tabularx}
\usepackage{multirow}
\usepackage{enumitem}
\usepackage{xcolor}
\usepackage{pgfplots}
\pgfplotsset{compat=1.18}

\newcolumntype{P}[1]{>{\RaggedRight\arraybackslash}p{#1}}
\newcolumntype{C}[1]{>{\centering\arraybackslash}p{#1}}

\sloppy

\title{Marketplace de Agentes Acadêmicos via MCP: Construção de um\\
Agente Acadêmico de Conteúdo para Cálculo II (Trilha B)}

\author{Nathã Barbosa de Souza Pinto, Cristiano Cardoso Cavalcante Lima}

\address{
  Instituto de Computação -- Universidade Federal do Amazonas (UFAM)\\
  Manaus -- AM -- Brasil\\
  \email{\{natha.pinto,cristiano.lima\}@icomp.ufam.edu.br}
  \\[6pt]
  \begin{minipage}{0.85\textwidth}
    \centering\footnotesize
    ICC220 -- Tópicos Especiais em Bancos de Dados / PPGINF528 -- Tópicos Especiais em
    Recuperação de Informação -- 2026/01\\
    Professores: André Luiz da Costa Carvalho e Altigran Soares da Silva
  \end{minipage}
}

\begin{document}

\maketitle

\begin{resumo}
Este trabalho descreve a construção de um Agente Acadêmico de Conteúdo (Trilha B) para a
disciplina Cálculo~II (IEM021) da UFAM, componente de um \textit{marketplace} de agentes
acadêmicos interoperáveis via \textit{Model Context Protocol} (MCP). O componente é formado por um
construtor de \textit{corpus} \textit{offline} e um servidor MCP \textit{read-only}, alcançando 100\%
de cobertura da ementa (7 tópicos, 304 \textit{chunks}, credibilidade média de 0,9398). Na Fase~2,
o servidor foi integrado a dois tutores reais sorteados e a um tutor extra voluntário, com 100\%
de sucesso técnico nas chamadas MCP completadas, revelando limitações de estabilidade temporal do
\textit{corpus} e uma falha de composição de resposta no lado do tutor.
\end{resumo}

\begin{abstract}
This paper describes the construction of an Academic Content Agent (Track B) for the Calculus~II
course (IEM021) at UFAM, a component of a marketplace of interoperable academic agents via the
Model Context Protocol (MCP). The component comprises an offline corpus builder and a read-only
MCP server, achieving 100\% syllabus coverage (7 topics, 304 chunks, average credibility of
0.9398). In Phase~2, the server was integrated with two drawn real tutors and one extra volunteer
tutor, with 100\% technical success on completed MCP calls, revealing temporal-stability
limitations of the corpus and a response-composition failure on the tutor side.
\end{abstract}

\section{Introdução}

Sistemas tutores automatizados enfrentam dois desafios centrais de Processamento de Linguagem
Natural: a coleta e organização de material de estudo a partir de fontes heterogêneas da
\textit{web} e a condução pedagógica fundamentada nesse material. Soluções baseadas exclusivamente
em \textit{Large Language Models} (LLMs) de propósito geral tendem a alucinar conteúdo de domínio,
não citam fontes verificáveis e não permitem inspecionar o conhecimento subjacente às respostas
\cite{mcp}.

Este relatório documenta o desenvolvimento de um \textbf{Agente Acadêmico de Conteúdo}
(Trilha~B, componente \texttt{B-G1}) para a disciplina \textbf{Cálculo~II (IEM021)} da UFAM,
integrado ao \textit{Marketplace de Agentes Acadêmicos via MCP} proposto nas disciplinas
ICC220/PPGINF528 (2026/01). O escopo entregue cobre as duas peças exigidas pelo enunciado: um
construtor de \textit{corpus offline} (\textit{parser} de ementa, \textit{planner}, coletor,
avaliador de fontes, indexador e congelador) e um servidor MCP \textit{read-only} que expõe o
\textit{corpus} congelado via o contrato comum da Seção~3.1 do enunciado. Cálculo~II foi escolhida
por três razões: (i) abundância de material acadêmico aberto de qualidade (OpenStax, MIT OCW,
Wikipedia); (ii) provas ENADE com gabarito oficial disponíveis pelo INEP, nas edições
\textbf{2021 (Matemática -- Licenciatura)} e \textbf{2014 (Matemática -- Bacharelado)}; e
(iii) tópicos objetivos, favoráveis à avaliação automática de acurácia. A combinação das duas
edições ENADE cobre tanto fundamentos (Licenciatura~2021) quanto os tópicos avançados da ementa
--- integrais múltiplas, Teorema de Green, séries de Taylor --- mais presentes na prova de
Bacharelado~2014.

\section{Metodologia}

\subsection{Disciplina-Alvo e Ementa}

IEM021 -- Cálculo~II é oferecida pelo Departamento de Matemática do ICE/UFAM (90h, obrigatória em
oito cursos). A ementa oficial organiza o conteúdo em sete tópicos: (I) Derivação de Vetores e
Regra da Cadeia; (II) Funções de Várias Variáveis; (III) Funções Potenciais e Integrais de Linha;
(IV) Derivadas de Ordem Superior; (V) Máximos e Mínimos; (VI) Integrais Múltiplas e Mudança de
Variáveis; (VII) Fórmula de Taylor e Séries.

\subsection{Arquitetura do Componente}

O \textbf{construtor de \textit{corpus}} é um \textit{pipeline} de sete módulos Python
(\texttt{python -m builder}): \texttt{parse\_ementa.py} (ementa~$\rightarrow$~JSON estruturado
com 31 subtópicos, conceitos, pré-requisitos e bibliografia); \texttt{planner.py} (plano de coleta
por tópico, \textit{queries} em PT/EN); \texttt{collector.py} (seis fontes, respeitando
\texttt{robots.txt} e \textit{rate limits}); \texttt{evaluator.py} (heurística de cinco
critérios); \texttt{dedup.py} (\textit{hash} exato e similaridade de \textit{shingles}, Jaccard
$\geq 0{,}85$); \texttt{indexer.py} (\textit{chunks} de 512 \textit{tokens}, sobreposição de 64,
ChromaDB, modelo \texttt{paraphrase-multilingual-MiniLM-L12-v2}); e \texttt{freeze.py}
(\texttt{corpus\_hash} SHA-256 determinístico). O valor padrão de recuperação $k=5$
(configurável até 20) segue literatura de RAG para tutoria.

O \textbf{servidor MCP de conteúdo} foi implementado com FastMCP \cite{fastmcp} sobre transporte
\textit{stdio}, com carregamento \textit{lazy} do ChromaDB para resposta imediata ao
\textit{handshake}. É estritamente \textit{read-only}: não expõe \textit{tools} de busca, coleta
ou escrita.

\begin{table}[h]
\centering
\caption{Tools do contrato MCP implementadas}
\label{tab:tools}
\small
\begin{tabularx}{\linewidth}{l l X}
\toprule
\textbf{Tool} & \textbf{Assinatura} & \textbf{Descrição} \\
\midrule
\texttt{list\_topics} & \texttt{()} & Tópicos, cobertura, credibilidade média, \texttt{corpus\_hash} \\
\texttt{corpus\_query} & \texttt{(query,k,filters)} & Busca semântica por cosseno; $k$ \textit{chunks} com metadados completos \\
\texttt{get\_chunk} & \texttt{(chunk\_id)} & Recupera \textit{chunk} específico por SHA-256 \\
\bottomrule
\end{tabularx}
\end{table}

\subsection{Avaliador de Fontes}

A heurística combina cinco critérios ponderados: domínio da fonte (0,30 -- OpenStax/MIT~OCW$=1{,}00$,
arXiv$=0{,}90$, Wikipedia$=0{,}85$); indícios de revisão editorial ou colaborativa (0,25); idioma
(0,15 -- PT$=1{,}00$, EN$=0{,}75$); presença de autor identificável (0,15); e volume de conteúdo
(0,15). Documentos com \textit{score}${}<0{,}40$ são reprovados; dos 26 documentos coletados,
todos foram aprovados, com \textit{score} médio de \textbf{0,9398}.

\section{Resultados da Fase 1}

O servidor MCP implementa corretamente as três \textit{tools} do contrato e os metadados
obrigatórios por \textit{chunk} (URL, \textit{score}, data, tópico, \textit{hash}, disciplina e
\texttt{corpus\_hash}). Testes unitários confirmaram conformidade com os requisitos R01--R12 do
\texttt{check\_contract} (\texttt{evaluation/contract/report.json}).

Com $N=3$ documentos mínimos por tópico, a cobertura da ementa atingiu \textbf{100\%} (7/7), com
304 \textit{chunks} indexados em 66,58~s de execução (Tabela~\ref{tab:cobertura}).

\begin{table}[h]
\centering
\caption{Cobertura, diversidade e credibilidade por tópico}
\label{tab:cobertura}
\small
\begin{tabularx}{\linewidth}{X c c}
\toprule
\textbf{Tópico} & \textbf{Docs} & \textbf{Cred. média} \\
\midrule
I -- Derivação de Vetores e Regra da Cadeia & 20 & 0,9089 \\
II -- Funções de Várias Variáveis & 85 & 0,9460 \\
III -- Funções Potenciais e Integrais de Linha & 60 & 0,9488 \\
IV -- Derivadas de Ordem Superior & 20 & 0,9074 \\
V -- Máximos e Mínimos & 28 & 0,9545 \\
VI -- Integrais Múltiplas e Mudança de Variáveis & 46 & 0,9522 \\
VII -- Fórmula de Taylor e Séries & 56 & 0,9606 \\
\bottomrule
\end{tabularx}
\end{table}

O gargalo do coletor foi o \textit{rate limit} da Wikipedia (429): 15/20 sucessos em PT e 0/11 em
EN, compensados por OpenStax (7/9), MIT~OCW (4/4) e arXiv (4/4); Khan Academy (0/4) falhou por
conteúdo renderizado via \textit{client-side JS}.

\section{Resultados da Fase 2}

O servidor \texttt{B-G1} foi disponibilizado a dois tutores sorteados -- \textbf{A-P1} (Giovanni
Escóssio Zanardo) e \textbf{A-G1} (Gabriel Dias Gonçalves e José Santarem) -- e a um tutor extra
voluntário, \textbf{A-G2}, fora do sorteio (bônus técnico da Seção~5 do enunciado; não substitui os
dois pares mínimos exigidos).

Cada parceiro regenerou o \textit{corpus} localmente (\texttt{python~-m~builder}), comportamento
esperado pois o índice ChromaDB não é versionado (Seção~6.4 do enunciado). Como o \textit{rate
limiting} da Wikipedia varia entre execuções, os \texttt{corpus\_hash} resultantes diferem entre
B-G1 (referência), A-P1 e A-G2, embora a estrutura (7 tópicos, mesmas fontes-alvo) permaneça
idêntica -- limitação discutida na Seção~\ref{sec:discussao}.

\subsection{Par A-P1 -- integração bem-sucedida}

O tutor A-P1 subiu o servidor localmente e aplicou a Suíte Estrutural (11 \textit{queries}) e a
Suíte de Conteúdo ENADE (20 questões: 6 da edição 2021--Licenciatura e 14 da edição
2014--Bacharelado, curadoria em \texttt{evaluation/content/enade\_curadoria\_calculo2.json}).
Todas as 62 chamadas retornaram sucesso, com 5 \textit{chunks} por \texttt{corpus\_query}.

\subsection{Par A-G1 -- falha técnica registrada}

O tutor A-G1 completou a coleta (25 documentos, 91,36~s) mas não concluiu o \textit{handshake} MCP
(\texttt{initialize()} travava indefinidamente). A causa mais provável é o \textit{timeout} do
cliente MCP de A-G1 sendo excedido pelo carregamento inicial dos \textit{embeddings} (471~MB) --
mesma causa do Erro~5 (Seção~\ref{sec:erros}). Como A-P1 conectou-se ao mesmo servidor e versão
sem problemas, descarta-se falha no servidor B-G1 como causa primária. O log da tentativa está em
\texttt{evaluation/interop/ag1\_collection\_log.json}.

\subsection{Par extra A-G2 (bônus, fora do sorteio)}

Tutor A-G2 (\texttt{qwen2.5-7b-instruct-mlx} via LM~Studio, temperatura~0), conectado via
\textit{stdio}. A auditoria de \textit{tools} confirmou uso restrito às três \textit{tools} do
contrato, sem chamadas não permitidas.

\begin{table}[h]
\centering
\caption{Chamadas MCP por parceiro real (tools comuns)}
\label{tab:interop}
\small
\begin{tabularx}{\linewidth}{l c c c c}
\toprule
\textbf{Parceiro} & \textbf{Tool} & \textbf{Chamadas} & \textbf{Sucesso} & \textbf{Lat. média} \\
\midrule
A-P1 & \texttt{list\_topics}  & 31 & 100\% & 10.274~ms \\
A-P1 & \texttt{corpus\_query} & 31 & 100\% & 28~ms \\
A-G2 & \texttt{list\_topics}  & 10 & 100\% & 10.316~ms \\
A-G2 & \texttt{corpus\_query} & 30 & 100\% & 58~ms \\
A-G2 & \texttt{get\_chunk}    & 22 & 100\% & 2,03~ms \\
\bottomrule
\end{tabularx}
\end{table}

\begin{figure}[h]
\centering
\begin{tikzpicture}
\begin{axis}[
    ybar, bar width=16pt, ymode=log,
    symbolic x coords={list\_topics,corpus\_query},
    xtick=data, width=0.85\linewidth, height=4.6cm,
    ylabel={Latência média (ms, escala log)},
    enlarge x limits=0.4,
    legend style={at={(0.5,1.18)},anchor=north,legend columns=-1,font=\footnotesize},
    nodes near coords, every node near coord/.append style={font=\tiny},
    ymin=1, ymax=20000,
]
\addplot coordinates {(list\_topics,10274) (corpus\_query,28)};
\addplot coordinates {(list\_topics,10316) (corpus\_query,58)};
\legend{A-P1, A-G2 (extra)}
\end{axis}
\end{tikzpicture}
\caption{Latência média por \textit{tool} nos dois parceiros com integração completa. O primeiro
\texttt{list\_topics} é dominado pelo carregamento único do modelo de \textit{embeddings}
(\textit{lazy loading}); \texttt{corpus\_query} opera na casa de dezenas de milissegundos em
ambos.}
\label{fig:latencia}
\end{figure}

A latência elevada do primeiro \texttt{list\_topics} ($\sim$10,3~s, Figura~\ref{fig:latencia})
reflete o custo único de carregar o modelo de \textit{embeddings} (471~MB); pago uma vez por
sessão graças ao \textit{lazy loading}, e reproduzido de forma consistente em dois ambientes e
modelos de tutor distintos (A-P1 e A-G2, este com \texttt{qwen2.5-7b-instruct-mlx}).

Na Suíte Estrutural aplicada por A-G2, 9/10 itens foram aprovados; o item~10 (recusa de orientação
médica) reprovou por citação indevida, discutida como Erro~6. A auditoria de citações registrou 15
alertas de baixo \textit{overlap} lexical entre afirmações do tutor e os \textit{chunks} citados,
concentrados em material com fórmulas matemáticas mal preservadas (Erro~4) -- evidência de que
aprovação de contrato mede interoperabilidade de protocolo, não qualidade pedagógica final.

\section{Análise de Erros}
\label{sec:erros}

\textbf{E1 -- \textit{Rate limit} Wikipedia (HTTP~429).} IPs de ambiente compartilhado (Colab)
ativaram bloqueios; menor diversidade PT nos tópicos~I e~VII; mitigado com \textit{delay} de
1,5~s e \textit{User-Agent} próprio.
\textbf{E2 -- URLs descontinuadas no OpenStax (404).} Reorganização entre edições reduziu
cobertura do tópico~I, compensada por MIT~OCW.
\textbf{E3 -- Khan Academy bloqueando \textit{scraping}.} Conteúdo renderizado via React;
solução futura: Selenium/Playwright.
\textbf{E4 -- Símbolos matemáticos corrompidos.} \textit{Encoding} incorreto no BeautifulSoup
(\texttt{"â f â y"} em vez de $\partial f/\partial y$); mitigado forçando
\texttt{resp.encoding=resp.apparent\_encoding}.
\textbf{E5 -- Latência do primeiro \texttt{list\_topics}.} Carregamento do modelo de
\textit{embeddings} (471~MB); custo pago uma única vez por sessão via \textit{lazy loading} --
mesma causa da falha de A-G1.
\textbf{E6 -- Citação indevida em resposta de recusa (A-G2).} No item de recusa de orientação
médica, o tutor recusou corretamente mas anexou três citações do \textit{corpus} de Cálculo~II sem
relação com a alegação (\textit{overlap} lexical nulo). É falha de \textbf{geração no lado da
Trilha~A}: o servidor retornou \textit{chunks} tecnicamente válidos para a consulta recebida; o
defeito está na composição da resposta final, não na recuperação -- distinção relevante para a
Seção~\ref{sec:discussao}.

\section{Discussão}
\label{sec:discussao}

\textbf{Cobertura vs. tempo.} O \textit{pipeline} completou 100\% da cobertura em 66,58~s; o
gargalo foi o \textit{rate limit} da Wikipedia, compensado por OpenStax/MIT~OCW.
\textbf{Diversidade vs. credibilidade.} Tópicos com menor diversidade de domínios (I e VII, 2
fontes) mantiveram credibilidade${}>0{,}90$, pois a priorização de OpenStax/MIT~OCW garante
material revisado editorialmente.
\textbf{Agnosticismo de domínio.} Testes com consultas fora do domínio principal retornaram
\textit{chunks} semanticamente próximos sem erros, validando o design \textit{read-only}.
\textbf{Estabilidade temporal.} URLs do OpenStax/MIT~OCW podem mudar entre edições e o
\textit{rate limit} da Wikipedia varia por execução, tornando o \textit{corpus} parcialmente
não determinístico -- observado nos três \texttt{corpus\_hash} distintos da Fase~2. A estrutura
(7 tópicos, mesmas fontes-alvo) permanece estável; varia apenas o subconjunto de artigos da
Wikipedia efetivamente coletado.
\textbf{Contaminação de dados.} O ENADE é público e provavelmente integra corpora de pré-treino
de LLMs; por isso o \textit{corpus} de referência é fixado por \texttt{corpus\_hash}, \textit{seed}
42 e decodificação determinística (\textit{greedy}, temperatura~0), tornando a avaliação
rastreável mesmo sob esse risco.
\textbf{Licenciamento e proveniência.} As fontes priorizadas têm regimes distintos: OpenStax é
CC~BY (menor risco jurídico); MIT~OCW é CC~BY-NC-SA (restringe uso comercial); Wikipedia é
CC~BY-SA (exige mesma licença em obras derivadas); arXiv mantém direitos com autores/periódico. O
repositório versiona apenas metadados e \textit{chunks} tratados, sem republicar documentos
brutos (Seção~6.4 do enunciado), mitigando risco de violação de licença.
\textbf{Potencial de generalização.} Os três experimentos (A-P1 bem-sucedido, A-G1 com falha
técnica documentada, A-G2 extra) indicam que o contrato MCP foi implementado de forma desacoplada
da identidade do parceiro -- A-P1 e A-G2 usam \textit{stacks} de LLM totalmente distintas e
obtiveram 100\% de sucesso técnico nas chamadas completadas. A generalização observada é, porém,
parcial: concentra-se em uma mesma área (Matemática/Cálculo), favorecendo sobreposição vocabular.
O Erro~6 mostra que interoperabilidade de protocolo não garante qualidade pedagógica -- teste
futuro natural é conectar o servidor a um tutor calibrado para área distinta (e.g. História), onde
a ausência de \textit{chunks} relevantes deveria acionar admissão de incerteza em vez de
recuperação forçada.

\textbf{Limitações.} (i)~ausência de fontes em PT de alta qualidade (SciELO); (ii)~\textit{chunks}
do MIT~OCW com baixa densidade informacional (páginas de navegação); (iii)~símbolos matemáticos
convertidos imperfeitamente de MathML/\LaTeX; (iv)~ausência de teste de interoperabilidade com
tutor de área não relacionada a Matemática, limitando conclusões sobre agnosticismo de domínio a
disciplinas tematicamente próximas.

\section{Conclusão}

Este trabalho demonstrou a viabilidade de um Agente Acadêmico de Conteúdo interoperável para
Cálculo~II, com 100\% de cobertura dos sete tópicos da ementa, credibilidade média de 0,9398 e 304
\textit{chunks} indexados. A separação entre construtor de \textit{corpus} e servidor MCP garantiu
reprodutibilidade e estabilidade das respostas. Os experimentos de Fase~2 -- A-P1 (sucesso), A-G1
(falha técnica documentada) e A-G2 (extra) -- validaram o design \textit{read-only} e o contrato
MCP em múltiplos ambientes e modelos de tutor, revelando também que interoperabilidade de
protocolo não equivale a qualidade pedagógica garantida (Erro~6). Os principais desafios técnicos
-- \textit{rate limiting} de APIs públicas, fragilidade de URLs e custo de inicialização de
\textit{embeddings} -- foram documentados e mitigados, constituindo contribuição prática para
trabalhos futuros em recuperação de informação acadêmica.

\begin{thebibliography}{9}

\bibitem{openstax_calc2}
STRANG, G. \textbf{Calculus Volume 2}. OpenStax, 2016. Disponível em:
\url{https://openstax.org/books/calculus-volume-2}. Acesso em: 5 jul. 2026.

\bibitem{openstax_calc3}
STRANG, G. \textbf{Calculus Volume 3}. OpenStax, 2016. Disponível em:
\url{https://openstax.org/books/calculus-volume-3}. Acesso em: 5 jul. 2026.

\bibitem{mitocw}
MIT OPENCOURSEWARE. \textbf{18.02SC Multivariable Calculus}. MIT, 2010. Disponível em:
\url{https://ocw.mit.edu/courses/18-02sc-multivariable-calculus-fall-2010}. Acesso em: 5 jul. 2026.

\bibitem{chromadb}
CHROMA. \textbf{ChromaDB: the open-source embedding database}. 2023. Disponível em:
\url{https://www.trychroma.com}. Acesso em: 5 jul. 2026.

\bibitem{sentencetransformers}
REIMERS, N.; GUREVYCH, I. Sentence-BERT: sentence embeddings using siamese BERT-networks. In:
\textbf{EMNLP 2019}. p.~3982--3992. Disponível em: \url{https://arxiv.org/abs/1908.10084}.

\bibitem{mcp}
ANTHROPIC. \textbf{Model Context Protocol (MCP)}. 2024. Disponível em:
\url{https://modelcontextprotocol.io}. Acesso em: 5 jul. 2026.

\bibitem{fastmcp}
FASTMCP. \textbf{FastMCP: fast, Pythonic MCP server framework}. 2024. Disponível em:
\url{https://github.com/jlowin/fastmcp}. Acesso em: 5 jul. 2026.

\bibitem{inep2021}
INEP. \textbf{ENADE 2021 -- Matemática Licenciatura: prova e gabarito}. Disponível em:
\url{https://download.inep.gov.br/enade/provas_e_gabaritos/2021_PV_licenciatura_matematica.pdf}.
Acesso em: 5 jul. 2026.

\bibitem{inep2014}
INEP. \textbf{ENADE 2014 -- Matemática Bacharelado: prova e gabarito}. Disponível em:
\url{https://download.inep.gov.br/educacao_superior/enade/provas/2014/33_matematica_bacharelado.pdf}.
Acesso em: 5 jul. 2026.

\end{thebibliography}

\end{document}